\section*{第 \ref{ch:g11:redox} 章——氧化与还原}

\begin{solution}{exo:g11:redox:1}
\omterm{def:g11:redox:oxidant}{氧化剂}：能够得到\omterm{def:g9:inside-the-atom:nucleus}{电子}的物种。\omterm{def:g11:redox:oxidant}{还原剂}：能够
失去\omterm{def:g9:inside-the-atom:nucleus}{电子}的物种。氧化：失去\omterm{def:g9:inside-the-atom:nucleus}{电子}的过程。还原：得到
\omterm{def:g9:inside-the-atom:nucleus}{电子}的过程。
\end{solution}

\begin{solution}{exo:g11:redox:2}
\ce{Fe^{2+}}/\ce{Fe}; \ce{Ag+}/\ce{Ag}; \ce{I2}/\ce{I-};
\ce{Fe^{3+}}/\ce{Fe^{2+}}.
\end{solution}

\begin{solution}{exo:g11:redox:3}
\ce{Ag+ + e- <=> Ag}; \ce{Al^{3+} + 3e- <=> Al};
\ce{Fe^{3+} + e- <=> Fe^{2+}}; \ce{Cl2 + 2e- <=> 2Cl-}.
\end{solution}

\begin{solution}{exo:g11:redox:4}
发生氧化反应（失去\omterm{def:g9:inside-the-atom:nucleus}{电子}）；铁作\omterm{def:g11:redox:oxidant}{还原剂}。
\end{solution}

\begin{solution}{exo:g11:redox:5}
锌被氧化，是\omterm{def:g11:redox:oxidant}{还原剂}；\ce{Cu^{2+}} 被还原，是
\omterm{def:g11:redox:oxidant}{氧化剂}。每个锌\omterm{def:g7:atoms-and-molecules:atom}{原子}失去 2 个\omterm{def:g9:inside-the-atom:nucleus}{电子}。
\end{solution}

\begin{solution}{exo:g11:redox:6}
从 \ce{H2O2 -> H2O} 出发：% ce-unbalanced-ok (step 1 of the method)
氧：左边有 2 个 O，故右边加 \ce{2H2O}；氢：左边
2 个，右边 4 个，故左边加 \ce{2H+}；电荷：左边为 $+2$，
右边为 0，故左边加两个\omterm{def:g9:inside-the-atom:nucleus}{电子}：
\ce{H2O2 + 2H+ + 2e- <=> 2H2O}.
\end{solution}

\begin{solution}{exo:g11:redox:7}
\ce{I2 + 2e- -> 2I-} 和 \ce{2S2O3^{2-} -> S4O6^{2-} + 2e-}；两
者各转移 2 个\omterm{def:g9:inside-the-atom:nucleus}{电子}，故 \ce{I2 + 2S2O3^{2-} -> 2I- + S4O6^{2-}}。
\end{solution}

\begin{solution}{exo:g11:redox:8}
\ce{Cu -> Cu^{2+} + 2e-} 和 $2\times$（\ce{Ag+ + e- -> Ag}）：
\ce{Cu + 2Ag+ -> Cu^{2+} + 2Ag}。生成的 \ce{Cu^{2+}} \omterm{def:g9:ions:ion}{离子}呈蓝色。
\end{solution}

\begin{solution}{exo:g11:redox:9}
\ce{C6H8O6 -> C6H6O6 + 2H+ + 2e-} 和 \ce{I2 + 2e- -> 2I-}：
\ce{C6H8O6 + I2 -> C6H6O6 + 2I- + 2H+}。维生素 C 是\omterm{def:g11:redox:oxidant}{还原剂}。
\end{solution}

\begin{solution}{exo:g11:redox:10}
\ce{Al -> Al^{3+} + 3e-} 给出 3 个\omterm{def:g9:inside-the-atom:nucleus}{电子}，\ce{Cu^{2+} + 2e- -> Cu}
需要 2 个；最小公倍数为 6，故第一式乘以
2，第二式乘以 3：\ce{2Al + 3Cu^{2+} -> 2Al^{3+} + 3Cu}。
\end{solution}

\begin{solution}{exo:g11:redox:11}
\ce{H2O2 + 2H+ + 2e- -> 2H2O} 和 $2\times$（\ce{Fe^{2+} -> Fe^{3+} + e-}）：
\ce{H2O2 + 2H+ + 2Fe^{2+} -> 2H2O + 2Fe^{3+}}.
\end{solution}

\begin{solution}{exo:g11:redox:12}
$n(\ce{Zn}) = 1.30/65.4 = \qty{1.99e-2}{mol}$。每氧化一个锌\omterm{def:g7:atoms-and-molecules:atom}{原子}，
便沉积一个铜\omterm{def:g7:atoms-and-molecules:atom}{原子}，故 $n(\ce{Cu}) = \qty{1.99e-2}{mol}$，
且 $m(\ce{Cu}) = 1.99\times 10^{-2} \times 63.5 = \qty{1.26}{g}$。
\end{solution}

\begin{solution}{exo:g11:redox:13}
\ce{Cr2O7^{2-} + 14H+ + 6e- -> 2Cr^{3+} + 7H2O} 和
$6\times$（\ce{Fe^{2+} -> Fe^{3+} + e-}）：
\ce{Cr2O7^{2-} + 14H+ + 6Fe^{2+} -> 2Cr^{3+} + 6Fe^{3+} + 7H2O}.
每个重铬酸根\omterm{def:g9:ions:ion}{离子}对应 6 个 \ce{Fe^{2+}}：$6 \times 0.010 = \qty{0.060}{mol}$。
\end{solution}

\begin{solution}{exo:g11:redox:14}
\ce{2ClO- + 4H+ + 2e- -> Cl2 + 2H2O} 和 \ce{2Cl- -> Cl2 + 2e-}；相加
后除以 2，得 \ce{ClO- + Cl- + 2H+ -> Cl2 + H2O}。酸提供
\ce{H+} \omterm{def:g9:ions:ion}{离子}，而反应释放出氯气，一种有毒气体：
因此切勿将两种产品混用。
\end{solution}

\begin{solution}{exo:g11:redox:15}
\omterm{def:g9:inside-the-atom:nucleus}{电子}不能在溶液中自由移动
（见\cref{prop:g11:redox:no-free-electrons}）：锌\omterm{def:g7:atoms-and-molecules:atom}{原子}失去的\omterm{def:g9:inside-the-atom:nucleus}{电子}只能
被与\omterm{def:g9:periodic-table-first-look:metal}{金属}接触的 \ce{Cu^{2+}} \omterm{def:g9:ions:ion}{离子}获得。铜\omterm{def:g7:atoms-and-molecules:atom}{原子}
因此在锌片表面生成，并留在那里。
\end{solution}

\begin{solution}{pb:g11:redox:1}
\textbf{1.} \ce{Cr2O7^{2-}}/\ce{Cr^{3+}} 和 \ce{CH3COOH}/\ce{CH3CH2OH}。

\textbf{2.} 乙醇被氧化：它是\omterm{def:g11:redox:oxidant}{还原剂}。

\textbf{3.} \ce{Cr2O7^{2-} + 14H+ + 6e- <=> 2Cr^{3+} + 7H2O}。

\textbf{4.} 碳\omterm{def:g7:atoms-and-molecules:atom}{原子}已配平（两边各 2 个）。氧：左边 2 个，
右边 1 个，故右边加 \ce{H2O}。氢：左边 4 个，
右边 $6 + 2 = 8$ 个，故左边加 \ce{4H+}。电荷：左边为 $+4$，
右边为 0，故左边加四个\omterm{def:g9:inside-the-atom:nucleus}{电子}：
\ce{CH3COOH + 4H+ + 4e- <=> CH3CH2OH + H2O}.

\textbf{5.} 共 12 个\omterm{def:g9:inside-the-atom:nucleus}{电子}：$2\times$ 第一式，$3\times$ 第二式
反向：
\ce{2Cr2O7^{2-} + 3CH3CH2OH + 16H+ -> 4Cr^{3+} + 3CH3COOH + 11H2O}.

\textbf{6.} 橙色的 \ce{Cr2O7^{2-}} \omterm{def:g9:ions:ion}{离子}被还原为绿色的
\ce{Cr^{3+}} \omterm{def:g9:ions:ion}{离子}，凡遇乙醇之处皆如此。

\textbf{7.} $M = 2\times 39.1 + 2\times 52.0 + 7\times 16.0 =
\qty{294.2}{g/mol}$.

\textbf{8.} $n = 2.0\times 10^{-3}/294.2 = \qty{6.8e-6}{mol}$.

\textbf{9.} 每 2 个重铬酸根\omterm{def:g9:ions:ion}{离子}对应 3 个乙醇\omterm{def:g7:atoms-and-molecules:molecule}{分子}：
$n = \tfrac{3}{2}\times 6.8\times 10^{-6} = \qty{1.0e-5}{mol}$.

\textbf{10.} $M(\ce{C2H6O}) = \qty{46.0}{g/mol}$，则
$m = 1.02\times 10^{-5}\times 46.0 = \qty{4.7e-4}{g}$，约为
\qty{0.47}{mg}。

\textbf{11.} $0.25/0.47 \approx 0.53$：约占重铬酸根的
\qty{53}{\percent}。

\textbf{12.} $0.53 \times 10 \approx \qty{5.3}{cm}$.

\textbf{13.} $\qty{0.25}{mg} \times 2100 = \qty{525}{mg}$，每升血液
中约为 \qty{0.53}{g/L}。

\textbf{14.} 呼出的气体首先在入口处接触颗粒；
该处重铬酸根过量，能将乙醇全部除去，因此反应
在管的前段即已完全进行，该段变为全绿，
然后乙醇才会到达更远处。

\textbf{15.} 反应是彻底的：绿色的 \ce{Cr^{3+}} \omterm{def:g9:ions:ion}{离子}不会再变回
重铬酸根，因此用过的管子不能再次测量。

% ledger: ghs:K2Cr2O7 (H350 among its hazard statements)
\textbf{16.} 该物质具有急性毒性（GHS06），并且是严重的健康危害
（GHS08：其危险说明包括“可能致癌”），此外还具有\omterm{def:g9:metals-acids-corrosion:corrosion}{腐蚀}性，
并对水生生物有毒；供公众使用的装置
不得含有该物质。

\textbf{17.} \ce{O2 + 4H+ + 4e- -> 2H2O} 和
\ce{CH3CH2OH + H2O -> CH3COOH + 4H+ + 4e-}：
\ce{CH3CH2OH + O2 -> CH3COOH + H2O}。

\textbf{18.} 4 个\omterm{def:g9:inside-the-atom:nucleus}{电子}。

\textbf{19.} 乙醇的\omterm{def:g10:the-mole:amount}{物质的量} $n = 0.25\times 10^{-3}/46.0 = \qty{5.4e-6}{mol}$；
\omterm{def:g9:inside-the-atom:nucleus}{电子}的\omterm{def:g10:the-mole:amount}{物质的量} $4n = \qty{2.2e-5}{mol}$，即
$2.2\times 10^{-5}\times 6.02\times 10^{23} = \num{1.3e19}$ 个\omterm{def:g9:inside-the-atom:nucleus}{电子}。

\textbf{20.} $Q = 1.3\times 10^{19}\times 1.60\times 10^{-19} \approx
\qty{2.1}{C}$。每个乙醇\omterm{def:g7:atoms-and-molecules:molecule}{分子}恰好通过导线输送 4 个\omterm{def:g9:inside-the-atom:nucleus}{电子}，
因此电荷量与乙醇的\omterm{def:g10:the-mole:amount}{物质的量}成正比，
测定电荷量即是测定乙醇。
\end{solution}
